\documentclass[11pt]{article}
\usepackage[margin=25mm]{geometry}
\usepackage{fontspec,unicode-math}
\setmainfont{Latin Modern Roman}
\setmathfont{Latin Modern Math}
\usepackage{amsmath,amssymb,mathtools}
\usepackage{xurl,hyperref}
\hypersetup{colorlinks=true,urlcolor=blue,linkcolor=blue}
\setcounter{secnumdepth}{0}
\setlength{\parindent}{0pt}
\setlength{\parskip}{5pt}
\setlength{\emergencystretch}{3em}
\newcommand{\sref}[2]{\hyperlink{source:#1:#2}{[S#2]}}
\newcommand{\cataloguescope}[1]{\paragraph{Source question.}#1}
\newcommand{\floor}[1]{\left\lfloor#1\right\rfloor}
\begin{document}
\section{Circuit-Size Hierarchy in Bounded Arithmetic}\label{statement}

\subsection{Definitions and mathematical statement}

\paragraph{Shared notation.}Standard number systems are denoted by $\mathbb N,\mathbb Z,\mathbb Q,\mathbb R,\mathbb C$; index ranges, computational models and specialized notation follow the cited problem source and the definitions below.


A Boolean circuit is an acyclic graph of Boolean gates computing a function on $n$ input bits; circuit size counts gates. A circuit-size hierarchy theorem separates classes of functions computable within different asymptotic circuit-size budgets, under the source gate and encoding conventions. The bounded-arithmetic theory $\mathrm{PV}_1$ formalizes polynomial-time reasoning with symbols for polynomial-time functions and its specified induction axioms. A language is in polynomial time if membership is decided by a uniform polynomial-time algorithm. An explicit circuit lower bound supplies such a uniform language together with a proved lower bound for its circuit sizes.
\paragraph{Statement recorded in the supplied source.}
Does the bounded-arithmetic theory $\mathrm{PV}_1$ prove the Circuit-Size Hierarchy theorem? If it does, can one extract from such a proof explicit polynomial-time languages with the corresponding circuit lower bounds? \sref{statement}{2}
\paragraph{Scope qualification.}
The formal hierarchy statement, size bounds and proof-extraction claim follow the cited source. A nonconstructive existence proof does not by itself produce a uniform polynomial-time hard language.
\paragraph{Source status.}
The supplied catalogue labels this record open. Its source-status metadata was checked against the supplied snapshot on 2026-10-11. This import does not independently establish current open status or mathematical correctness. The catalogue literature assessment is dated 2026-08-22.
\subsection{Short English statement}
Can polynomial-time bounded arithmetic prove a circuit-size hierarchy and yield explicit efficiently decidable functions realizing its lower bounds?
\subsection{Sources}
\begin{itemize}
\item[\hypertarget{source:statement:1}{S1}] Ulam AI and the UnsolvedMath Contributors, UnsolvedMath: A Curated Collection of Open Mathematics Problems, catalogue entry. Catalogue curation and metadata: CC BY 4.0. The supplied source question is retained; definitions and the English summary are editorial additions.
\url{https://huggingface.co/datasets/ulamai/UnsolvedMath}.
\item[\hypertarget{source:statement:2}{S2}] Proof Complexity and Beyond (2024), Oberwolfach Reports, DOI 10.4171/owr/2024/15, pp. 871--934
\url{https://doi.org/10.4171/owr/2024/15}.
\item[\hypertarget{source:statement:3}{S3}] Carmosino et al., Provability of the Circuit Size Hierarchy
\url{https://doi.org/10.4230/LIPIcs.ITCS.2025.30}.
\item[\hypertarget{source:statement:4}{S4}] Oberwolfach Report 15/2024
\url{https://oa.tib.eu/renate/bitstreams/b0ef46dc-1bdf-45b7-a4f5-ca0a9e60a7d6/download}.
\textit{Underlying Oberwolfach Reports material retains its source-specific licence, including CC BY-SA 4.0 where applicable; see THIRD\_PARTY\_NOTICES.md.}
\end{itemize}
\paragraph{Preparation and review.}
Original question attributed to the cited source and UnsolvedMath contributors. Codex (GPT-6; exact serving revision not exposed) assisted with definitions, English summaries, source inspection and formatting. No new human mathematical authorship or independent proof verification is asserted. Catalogue identity checked against the complete supplied registry and existing repository statements; known flagged repetitions and obvious identical conjectures excluded. Full mathematical-equivalence/literature review has not been performed. Statement preparation and final review followed contribution prompts C and D; automated checks verify structure and identity, not open status.
\label{end:statement}
\end{document}
